\section{Two temperature tunings}
\label{sec:tuning}

The scale estimates select a temperature at which the perturbation
decays.  For a general law \(\rho\), there is no reason for that selection
to be self-dual.  We therefore apply the selection separately to the
physical family and to its exact dual, using the same physical inverse
temperature as the parameter.  This section constructs the two
trajectories; their temperatures will be compared afterward.

\subsection{Primal and dual initial arrays}

Write \(K_e(\beta)=\beta J_e=\beta(1+\varepsilon\xi_e)\) for the
coupling including inverse temperature.  For \(K>0\), define its dual
coupling \(g(K)\) by
\begin{equation}
 e^{2g(K)}-1=\frac{2}{e^{2K}-1}.
 \label{eq:dual-coupling}
\end{equation}
Thus \(g\) is a smooth decreasing involution, \(g(K_0)=K_0\),
\(g'(K_0)=-1\), and \(g''(K_0)=2\sqrt2\).
The edge odds and their duals are
\[
 w_e(\beta)=e^{2K_e(\beta)}-1,
 \qquad w^*_{e^*}(\beta)=\frac2{w_e(\beta)},
 \qquad K^*_{e^*}(\beta)=g(K_e(\beta)),
\]
where \(e^*\) crosses \(e\).  We identify the dual square lattice with
a translate of the square lattice.

Use the scale-step parameters of Proposition~\ref{prop:scale-estimates}
and choose an integer starting spacing \(s_0\) sufficiently large for
the construction.  We also take it large enough for the fixed geometric
checks used later in the crossing estimates of \cite[Section~7]{R}.
Both lattices use the spacings
\(s_k=s_0L^k\).  The two input families, parametrized by a small
number \(h\), are
\begin{equation}
 \begin{array}{c|c|c}
  \text{family}&\text{physical inverse temperature}
       &\text{couplings supplied to the scale map}\\ \hline
  \mathrm p&K_0+h&(K_0+h)(1+\varepsilon\xi_e)\\
  \mathrm d&K_0-h&g((K_0-h)(1+\varepsilon\xi_e)).
 \end{array}
 \label{eq:two-input-families}
\end{equation}
The opposite sign of \(h\) in the second row makes the mean charge
increase with \(h\) in both initializations.  No symmetry of the
distribution is involved in this parametrization.

Write \(\widetilde K_e\) for the coupling in either row, on the lattice
of that row.  Allocate each edge to its nearest \(s_0\)-grid center,
with distance measured from its midpoint and ties split equally.  If
\(\alpha_{i,e}\) is its allocation to \(s_0i\), then
\(\alpha_{i,e}\ge0\), \(\sum_i\alpha_{i,e}=1\), and the allocation
respects grid translations and square symmetries.  Define
\begin{equation}
 f_i^{(0)}=\sum_e\alpha_{i,e}(\widetilde K_e-K_0)b_e.
 \label{eq:initial-array}
\end{equation}
The sum of these functions is precisely the change from the pure Ising
logarithmic weight, up to a spin-independent constant.  This remains
true for conditional specifications with exterior pins, including the
bonds incident to the free spins.

Each initial potential is bounded, even, pure centered, and determined
by bonds in a fixed multiple of its cell.  The iid bond law and the
equivariant allocation give grid stationarity and square symmetry in
law.  For the dual array, the crossing-edge correspondence is a bijection,
so the transformed bonds are again iid on the dual lattice.  All the
array hypotheses of Section~\ref{sec:scale} therefore hold for both
families.  They persist under iteration.  Moreover, boundedness of the
initial disorder and \eqref{eq:raw-scale-bound} give a deterministic
bound on each potential at every fixed number of steps, uniformly on a
compact interval of \(h\).

Recall \(\sigma_\rho^2=\E_\rho\xi^2>0\), and put
\[
 C_\rho=2K_0^2\sigma_\rho^2.
\]
The strict positivity is where nondegeneracy of \(\rho\) enters the
tuning.  The following estimates also locate both temperatures to
second order in the disorder strength.

\begin{lemma}[Initial coordinates]
\label{lem:initial-coordinates}
Fix \(T<\infty\).  Uniformly for \(|h|\le T\varepsilon^2\), either
family in \eqref{eq:two-input-families} satisfies
\begin{equation}
 \begin{aligned}
 X_0&\le C(s_0)\varepsilon+O_{T,s_0}(\varepsilon^2),
 &Y_0&=O_{T,s_0}(\varepsilon^2),\\
 m_0&=2s_0\bigl(h+d_\rho\varepsilon^2+O_T(\varepsilon^3)\bigr),
 &W_0&=C_\rho\varepsilon^2+O_T(\varepsilon^3),
 \end{aligned}
 \label{eq:initial-coordinate-bounds}
\end{equation}
where \(d_\rho=0\) for the primal family and
\(d_\rho=\tfrac12g''(K_0)K_0^2\sigma_\rho^2\) for the dual family.
For every subsequently fixed \(H>R_0\),
\begin{equation}
 v_{H,0}=C_\rho\varepsilon^2+O_{T,H,s_0}(\varepsilon^3).
 \label{eq:initial-variance}
\end{equation}
\end{lemma}

\begin{proof}
The primal coupling displacement is
\[
 \widetilde K_e-K_0=h+K_0\varepsilon\xi_e+O_T(\varepsilon^3).
\]
For the dual family, Taylor expansion at \(K_0\) gives
\begin{equation}
 \widetilde K_e-K_0
 =h-K_0\varepsilon\xi_e
   +\tfrac12g''(K_0)K_0^2\varepsilon^2\xi_e^2
   +O_T(\varepsilon^3).
 \label{eq:dual-initial-expansion}
\end{equation}
The remainders are uniform because \(|\xi_e|\le1\).  In particular,
the dual mean contains a quadratic shift even though \(\E_\rho\xi=0\).
There is no claim that its cubic remainder vanishes.

There are two edges per square-lattice vertex.  Averaging the allocation
over a period consequently gives
\(\sum_e\alpha_{0,e}=2s_0^2\).  The charge normalization and
independence of different edges give the exact identities
\begin{align}
 m_0&=2s_0\E(\widetilde K_e-K_0),\notag\\
 W_0&=s_0^{-2}\sum_e\alpha_{0,e}
          \left(\sum_i\alpha_{i,e}\right)\Var(\widetilde K_e)
      =2\Var(\widetilde K_e).
 \label{eq:initial-charge-variance}
\end{align}
Substitution of the two expansions proves the asserted bounds for
\(m_0,W_0\).  The finite number of uniformly bounded functions \(b_e\)
in a cell gives the bounds for \(X_0,Y_0\); the coefficient of the
leading \(\varepsilon\) bound for \(X_0\) is independent of \(T\).
Finally, \eqref{eq:variance-near-W} yields
\eqref{eq:initial-variance}.
\end{proof}

\subsection{Keeping the expanding coordinate between two walls}

We now adapt the shooting argument of \cite[Section~6]{R} to these
two initializations.  The purpose of the argument is to control the
expanding mean charge while the variance slowly decreases.  Its
continuity step must accommodate atoms in \(\rho\), and will use the
independent continuous thresholds of the scale construction.

\begin{proposition}[Two small trajectories]
\label{prop:two-trajectories}
Let \(\rho\) be a nondegenerate mean-zero law supported on \([-1,1]\).
There are constants \(A,B,D,c>0\), a fixed cutoff \(H\), and
\(\varepsilon_0(\rho)\in(0,1)\) with the following property.  For each
\(0<\varepsilon<\varepsilon_0(\rho)\), there are positive deterministic physical
inverse temperatures
\begin{equation}
 \beta_{\mathrm p}=K_0+O_\rho(\varepsilon^2),
 \qquad \beta_{\mathrm d}=K_0+O_\rho(\varepsilon^2)
 \label{eq:tuned-temperatures}
\end{equation}
such that the primal array at \(\beta_{\mathrm p}\) and the dual array
at \(\beta_{\mathrm d}\) have, at the common spacings \(s_k=s_0L^k\),
\begin{equation}
 v_k:=v_{H,k}>0,\qquad
 X_k\le A\sqrt{v_k},\qquad Y_k\le Bv_k,\qquad |m_k|\le Dv_k,
 \label{eq:trajectory-region}
\end{equation}
and
\begin{equation}
 v_k\le\frac1{v_0^{-1}+ck}.
 \label{eq:trajectory-decay}
\end{equation}
In particular, \(X_k\to0\) for both arrays.  Given any \(\eta>0\),
the same choices can be made with \(\sup_k X_k\le\eta\) by decreasing
\(\varepsilon_0(\rho)\).  All choices precede the specification of a
domain or its lattice approximations.
\end{proposition}

\begin{proof}
We arrange that the two norm bounds persist whenever \(|m|\le Dv\),
while the map sends the two mean walls strictly outward.  Continuity
will then select a parameter that never exits the region.  The constants
\(A,B\) control the norm bounds, \(D\) controls the mean walls, and
the cutoff \(H\) makes the variance decrease.

We choose constants common to the two families.  First take \(A\)
large enough for the leading initial bound \(X_0/\sqrt{v_{H,0}}\),
and so that
\begin{equation}
 \lambda A+2C_1<A.
 \label{eq:choose-A}
\end{equation}
This choice is possible before \(T,H\) are fixed: the leading variance
is \(C_\rho\varepsilon^2\), and the leading norm bound in
Lemma~\ref{lem:initial-coordinates} is independent of those parameters.
Next choose \(D>0\) so that
\begin{equation}
 (L-1)D>C_1A^2+1.
 \label{eq:choose-D}
\end{equation}
Choose \(T\) large enough that, in both families, the initial ratio
\(m_0/(Dv_{H,0})\) is below \(-1\) at \(h=-T\varepsilon^2\)
and above \(1\) at \(h=T\varepsilon^2\), for sufficiently small
\(\varepsilon\) and every subsequently fixed \(H\).  Indeed its
leading value is
\[
 \frac{2s_0}{DC_\rho}
       \left(\frac h{\varepsilon^2}+d_\rho\right),
\]
an affine function with positive slope.  Finally choose \(B\) large
enough for \(Y_0\le Bv_{H,0}\) on this full interval, and for
\begin{equation}
 \lambda B+C_1(D+A^2)<B.
 \label{eq:choose-B}
\end{equation}
The order of choices so far is \(A,D,T,B\).

Suppose an array satisfies the inequalities
\eqref{eq:trajectory-region}, with \(v=v_H\) sufficiently small.
Then \eqref{eq:variance-near-W} gives
\begin{equation}
 W=v+O_H(v^{3/2}).
 \label{eq:W-along-region}
\end{equation}
Apply \eqref{eq:marginal-variance-flow} with \(z=\sqrt v\).  Since
\[
 \sum_{H<|y|\le LH}|y|^{-2}\longrightarrow2\pi\log L>0,
\]
we may now choose \(H\) sufficiently large that the deterministic bound
on \(|e_H|\) is less than one quarter of
\(2\kappa\sum_{H<|y|\le LH}|y|^{-2}\).  After fixing this \(H\),
restrict \(v\) to a sufficiently
small interval, ensuring also \(A\sqrt v\le\min\{1,x_*\}\).
Equation~\eqref{eq:W-along-region} and the \(O_H(v^{5/2})\) remainder
then imply, for fixed \(c,C_2>0\),
\begin{equation}
 v-C_2v^2\le v'\le v-cv^2,
 \qquad v'>0,\qquad \frac{v'}v=1+O(v).
 \label{eq:variance-trap}
\end{equation}
Positivity follows by requiring \(C_2v<1/2\).  All these restrictions
are uniform in the scale.

The norm inequalities point into the region.  Indeed,
\[
 X'\le\lambda A\sqrt v+C_1\sqrt W+C_1Bv
     \le(\lambda A+C_1+o(1))\sqrt v
     \le A\sqrt{v'},
\]
by \eqref{eq:choose-A} and \(v'/v\to1\).  Similarly,
\[
 Y'\le\bigl(\lambda B+C_1(D+A^2)\bigr)v\le Bv'
\]
by the strict slack in \eqref{eq:choose-B}.  Conversely the two mean
walls point outward: at \(m=Dv\),
\[
 m'\ge LDv-C_1A^2v>Dv',
\]
and at \(m=-Dv\) the corresponding upper bound gives \(m'<-Dv'\).
Here \eqref{eq:choose-D} and \eqref{eq:variance-trap} supply the strict
inequalities.  Thus the variance and norm requirements persist up to
and including a first exit through either mean wall.

We have fixed \(A,D,T,B,H\) and a small interval of permitted \(v\).
Only now decrease \(\varepsilon_0(\rho)\).  The initial estimates then
give positive \(v_0\) in this interval, together with both norm bounds,
for every \(h\in[-T\varepsilon^2,T\varepsilon^2]\).  They also give
the two strict endpoint inequalities for \(m_0/(Dv_0)\).

To apply a shooting argument we next verify continuity in \(h\) at
every finite time.  Couple nearby parameters using the same bond
variables, construction priorities, and thresholds.  At a fixed number
of steps, a prescribed finite collection of determining stencils uses
a deterministic finite set of spin sites and potential coin labels.
The comparison graph is chosen independently of the flag statuses;
unused coins remain as dummy variables, and integrations retain all old
field variables.  Thus all the potentials in these stencils can be
viewed on a common finite field alphabet, independently of \(h\).
Fix a parameter \(h\).  Only finitely many comparisons determine these
potentials.
Priorities have no ties almost surely.  At each flag comparison, its
threshold is independent of the input data and has a continuous law,
so equality has probability zero, even when \(\rho\) has atoms.
Inductively, the statuses are locally constant at the fixed \(h\)
almost surely, and all remaining finite conditional sums and logarithms
are continuous.  Consequently the potentials are almost surely
continuous at this fixed \(h\) in supremum norm.  The deterministic
finite-step bound following \eqref{eq:initial-array} permits dominated
convergence.  In particular, as \(h_n\to h\),
\[
 \bigl\|\E f_0^{(k)}(h_n)-\E f_0^{(k)}(h)\bigr\|_\infty
 \le\E\bigl\|f_0^{(k)}(h_n)-f_0^{(k)}(h)\bigr\|_\infty
 \longrightarrow0,
\]
which proves continuity of \(Y_k\).  The same domination, the bounded
linear charge, and the finite covariance and coefficient sums prove
continuity of \(X_k,m_k,W_k,v_{H,k}\).

For either family, this continuity first gives a closed interval
\(I_0\subset[-T\varepsilon^2,T\varepsilon^2]\) on which
\(-1\le m_0/(Dv_0)\le1\), with the left and right endpoints on the
lower and upper walls, respectively.  Suppose a nonempty closed
interval \(I_k\) has been chosen so that every parameter in it satisfies
\eqref{eq:trajectory-region} through time \(k\), and its endpoints
lie on the opposite time-\(k\) walls in that order.  The estimates
above give \(v_{k+1}>0\) and the two norm bounds throughout \(I_k\).
Its endpoints have
\(m_{k+1}/(Dv_{k+1})<-1\) and \(>1\), respectively.  Take the first
point, starting at the left endpoint, where this continuous ratio is
\(1\), and then its last preceding point where the ratio is \(-1\).
Between those points the ratio lies in \([-1,1]\).  They therefore
bound a closed interval \(I_{k+1}\subset I_k\) with all the required
properties at time \(k+1\); the earlier-time requirements are retained.
The identical first-and-last-hit construction gives \(I_0\).

Compactness supplies \(h\in\bigcap_{k\ge0}I_k\).  No monotonicity of
the iterated coordinates in \(h\) is needed.  For this parameter,
\eqref{eq:variance-trap} holds at every step, and hence
\[
 \frac1{v_{k+1}}\ge\frac1{v_k}+c.
\]
This proves \eqref{eq:trajectory-decay} and \(X_k\to0\).
Since \(v_k\le v_0=O_\rho(\varepsilon^2)\), it also proves the
arbitrarily small uniform bound on \(X_k\).

Apply this argument to both rows of \eqref{eq:two-input-families},
obtaining parameters \(h_{\mathrm p}\) and \(h_{\mathrm d}\).
The physical temperatures are
\(\beta_{\mathrm p}=K_0+h_{\mathrm p}\) and
\(\beta_{\mathrm d}=K_0-h_{\mathrm d}\).  Their positivity and
\eqref{eq:tuned-temperatures} follow from
\(|h_{\mathrm p}|,|h_{\mathrm d}|\le T\varepsilon^2\).
\end{proof}

The shooting construction alone does not establish uniqueness.
Each selected temperature gives an exact
small-potential representation at every scale, with the perturbation
vanishing in environmental moments.  The next argument compares these
representations on a common torus and forces their physical temperatures
to agree.
